\scanpage{5}
(v) Si $X\in D_\ell(A_K)$ \uline{ou} $\underline a\in Z_\ell^{\unsure{*|'}}(A_K)$, on a
$\unsure{v|\nu}(X,\underline a)\in\mathbb Z$.

% slanted
\margin{\deleted{\unsure{pour} $X$, \illegible{1} $\underline a$ \unsure{fixés}, \illegible{1} \unsure{ne dépend}}}

(vi) Si $m={}\text{degré de }A_0/A_0^0$, alors on a
tjrs $m v(X,\underline a)\in\mathbb Z$.

[N.B. On obtient donc \inserted{par (v) et (vi)} un accouplement remarquable

\[
\underset{\substack{\shortparallel\\ A^*(K)}}{D_a(A_K)/D_\ell(A_K)}
 \times Z_0^*(A_K)/Z_\ell^*(A_K)\longrightarrow
 \underset{\substack{\shortparallel\\ {}_m(\mathbb Q/\mathbb Z)}}{\mathbb Z/m\mathbb Z}
\]

% slanted
\margin{À examiner : Cas d’un fibré principal homogène sous une VA \unsure{$\downarrow$|et}}

\deleted{pour \unsure{une} \unsure{valeur|unicité} \unsure{$0)$|$?)$}} par suite

\[
 A^*(K)\times A(K)\longrightarrow\mathbb Z/m\mathbb Z \qquad ]\tag*{$*$}
\]

La relation avec le symbole local $(X,\underline a)_v$ est
donnée par

\[
 \boxed{(X,\underline a)_v=v(X,\underline a)+\sum_{x_0\in A_0^{\unsure{*|c}}}i(\bar X\unsure{,|\,\cdot\,}\overline{\underline a},x_0)}\tag*{(vii)}
\]

% slanted
\margin{Se généralise aussitôt à des variétés quelconques \textendash\textendash\textendash}

Ceci montre que $(X,\underline a)_v\in\frac{1}{m}\mathbb Z$ pour tt $X,\underline a$,
\deleted{donc} à fortiori il est \uline{rationnel}, et il est
entier si $X\in D_\ell(A_K)$ [On le savait
déjà, car on trouve, si $X=\operatorname{div} f$, $(X,\underline a)_v=v(f(\underline a))$]
\uline{ou} si $\underline a\in\deleted{\unsure{D}}Z_\ell^*(A_K)$. On trouve donc
% archive sheet number: 4
